% ===================================================================
%  HOTO Na 2 — Jikin ɗan adam. --- Lokacin hutu. Wasanni.
%
%  Traduction, faite en 2026, en haoussa d'aujourd'hui, sur l'ido de
%  texto/io. Regles de la colonne : voir 10-hoto-01.tex. Une seule
%  numerotation.
%
%  LE CORPS EST HAOUSSA D'UN BOUT A L'AUTRE, ORGANES COMPRIS, ET
%  C'EST L'EXACT CONTRAIRE DU PERSAN. Au meme tableau, la colonne
%  persane partageait le corps en deux : ce qu'on montre du doigt
%  portait un nom iranien, ce qu'il faut ouvrir pour le voir portait
%  un nom arabe — جمجمه le crane, قلب le coeur, معده l'estomac —,
%  parce que la medecine du monde iranien s'est ECRITE en arabe. Le
%  haoussa est musulman depuis mille ans lui aussi, il doit a l'arabe
%  sa religion, son ecole et son calendrier — mais pas un os :
%
%    ƙoƙon kai   le crane      ƙwaƙwalwa  le cerveau
%    zuciya      le coeur      huhu       le poumon
%    tumbi       l'estomac     hanji      l'intestin
%    haƙarƙari   les cotes     ƙashi      l'os
%    jijiya      la veine      jini       le sang
%    tsoka       le muscle     fata       la peau
%
%  Aucun n'est arabe. La medecine ne s'est pas ecrite ici, elle s'est
%  dite — et ce qui se dit reste. « ƙoƙon kai », c'est la CALEBASSE
%  de la tete ; « idon sawu », la cheville, c'est l'oeil de
%  l'empreinte. Le corps se nomme par figures et ne doit rien a
%  personne.
%
%  MAIS LES SEPT JOURS DE LA SEMAINE SONT ARABES, TOUS LES SEPT, SANS
%  EXCEPTION : Litinin, Talata, Laraba, Alhamis, Jumma’a, Asabar,
%  Lahadi. Le tableau en emploie deux, le jeudi et le dimanche. Le
%  persan du tableau 3 comptait ses jours — یکشنبه, دوشنبه, un et deux
%  apres le samedi — et ne laissait a l'arabe que le vendredi de la
%  priere. Deux langues musulmanes, deux reponses opposees a la meme
%  question, et la difference est que la semaine persane existait
%  AVANT l'islam.
%
%  LES CINQ SENS SONT DEUX, ET C'EST LA TROISIEME REPONSE DU RELEVE A
%  CET ALINEA. Le persan formait cinq noms sur cinq verbes par un seul
%  suffixe — بینایی, شنوایی, بویایی, چشایی, بساوایی ; le javanais
%  faisait cinq formes en pa- ; le haoussa n'en a qu'un et demi. Un
%  seul verbe, « ji », entend, sent, goute ET touche : « jin sauti »,
%  « jin wari », « jin ɗanɗano », « jin taɓawa » — et il sert encore a
%  la douleur, « ina jin ciwo », que la note de l'alinea 17 traduit
%  mot pour mot. Seul « gani », voir, se tient a part. La colonne doit
%  pourtant en nommer cinq, puisque la planche en grave cinq : elle
%  ecrit gani, ji, sunsuna, ɗanɗano, taɓawa. Le vocabulaire suit la
%  gravure ; la langue, elle, dit autre chose.
%
%  ET LA BICYCLETTE (37) EST LE CINQUIEME MOT DU RELEVE POUR CETTE
%  MACHINE, ET LE SEUL QUI NE SOIT PAS NEUF. Le javanais dit pit, du
%  neerlandais fiets ; l'indonesien sepeda, de velocipede ; le
%  francais bicyclette ; le persan دوچرخه, la deux-roues, et il
%  decrit. Le haoussa dit « keke », qui est le mot de la machine a
%  manivelle — « keken ɗinki » est la machine a coudre, « keken saƙa »
%  le metier a tisser. Quatre langues ont emprunte ou fabrique ; une
%  seule avait deja le mot et l'a preté.
%
%  TROIS OBJETS QUE LE HAOUSSA DECRIT COMME LE PERSAN LES DECRIVAIT,
%  ET LE MEME MOT DESSOUS : les echasses (24) sont « sandunan ƙafa »,
%  les bois-pieds, ce que le persan disait چوب‌پا ; la serre (26) est
%  « gidan tsirrai », la maison des plantes, ce que le persan disait
%  گلخانه ; le cerf-volant (33) est « tsuntsun iska », l'oiseau de
%  vent, la ou le persan disait بادبادک, le petit vent-vent. Trois
%  langues sans parente qui, devant trois objets venus, prennent trois
%  fois le meme parti.
%
%  ET DEUX CHOSES QUI N'ONT PAS DE NOM ICI, PARCE QU'ELLES NE SONT
%  JAMAIS VENUES : la toupie (18) et l'ours (32). La colonne ecrit
%  « abin jujjuyawa », la chose qui tourne, et garde « beyar », qui est
%  un emprunt pour une bete qu'aucun Haoussa n'a vue. Un mot emprunte
%  pour un animal absent en dit plus long qu'un mot refuse.
%
%  DEUX REGLAGES DE kolonoj.py, ET LES DEUX SONT VENUS DU TEXTE, NON
%  D'UNE PREVISION.
%
%  1. UNE REGLE RESSERREE, PARCE QU'UN MOT REDOUBLE N'EST PAS LE MEME
%     MOT. Le N.-B. du dernier alinea invite le maitre a decrire
%     chaque jeu « daki-daki », en detail. Or « daki » etait releve,
%     puisque la chambre s'ecrit « ɗaki » et que la forme nue ne
%     valait rien. Elle valait quelque chose : liee a elle-meme par un
%     trait d'union, elle est un adverbe. La regle ne s'ouvre donc
%     plus sur un mot qui touche un trait d'union, des deux cotes.
%     C'est le « rumah sakit » javanais dans une autre langue — la
%     regle avait raison sur le MOT et tort sur la LOCUTION —, et le
%     tableau 1 porte la correction a la place ou la liste avait ete
%     annoncee trop sure.
%
%  2. UNE REGLE QUI TIENT UN CHOIX, ET NON UNE FAUTE. Le haoussa a
%     deux orthographes pour une seule lettre : le Niger ecrit « ƴan »,
%     le Nigeria « ’yan ». Les deux sont justes chez eux. Ce tableau a
%     ete ecrit d'abord a la nigeriane, et la regle a crie trois fois.
%     Elle avait raison : la colonne s'est donnee QUATRE CROCHETS au
%     tableau 1, et trois crochets plus une apostrophe n'est pas le
%     meme alphabet. Les trois « ’yan » sont donc devenus « ƴan », et
%     ce que la regle releve desormais n'est pas une faute mais un
%     MELANGE. C'est la premiere du releve dans ce cas.
% ===================================================================

%%K t02-apar-1 apar
\VUsaut{4.0mm}
\VUtitre{15.0pt}{60}{HOTO Na 2}
\VUsaut{2.0mm}
\VUtitre{11.4pt}{0}{Jikin ɗan adam. --- Lokacin hutu.}
\VUtitre{11.4pt}{0}{Wasanni.}

%%K t02-tit-0 sub
\VUtitre{13.2pt}{0}{Jikin ɗan adam.}
\VUsaut{2.4mm}
\VUcentre{10.2pt}{0}{\textit{(Darasi mai sauƙi na ilimin halitta.)}}
\VUsaut{3.0mm}

%%K t02-tit-1 sub pos=t02-01-1
\VUpk{10.2pt}{40}{Kai.}
\VUsaut{3.0mm}

%%K t02-01-1 p
1. --- Manyan sassan jikin ɗan adam su ne : \VUgras{kai,}
\VUgras{gaɓɓai} da \VUgras{hannaye da ƙafafu.}

%%K t02-02-1 p
2. --- Kai yana da sassa biyu : \VUgras{ƙoƙon kai,} wanda yake ɗauke
da \VUgras{ƙwaƙwalwa} kuma \VUgras{gashi} ya rufe shi, da
\VUgras{fuska.}

%%K t02-03-1 p
3. --- A zanenmu ba ma ganin ƙoƙon kai ko ƙwaƙwalwa ; sai dai muna
ganin manyan sassan fuska sarai.

%%K t02-04-1 p
4. --- Ga \VUgras{goshi} da farko, wanda yake nuna hankali mai ƙarfi
da tunani mai aiki kullum. \VUgras{Gefen goshi} yana da fili sosai.
\VUgras{Ido} yana da rai kuma a buɗe yake sarai. \VUgras{Gira} mai
kauri da \VUgras{gashin ido} masu tsawo suna kiyaye shi.

%%K t02-05-1 p
5. --- Ga kuma \VUgras{hanci} da \VUgras{ramukan hanci.} Hanci yana
taimaka mana mu sunsuna mu kuma numfasa. A kowane gefen hanci akwai
\VUgras{kunci,} can nesa kuma \VUgras{kunne.} Ƙasan fuska ya ƙunshi
\VUgras{baki} da \VUgras{haɓa.}

%%K t02-06-1 p
6. --- Ba ma ganinsu sosai, domin \VUgras{gashin baki} da
\VUgras{gemu} sun ɓoye mana su. Amma mun san cewa baki yana da
\VUgras{leɓɓa} biyu, \VUgras{muƙamuƙin sama} da \VUgras{muƙamuƙin
ƙasa} tare da \VUgras{haƙoransu,} da kuma \VUgras{harshe.} Da baki
muke magana muke kuma cin abinci. Da harshe muke jin ɗanɗanon abinci.

%%K t02-07-1 p
7. --- Ido, kunne, hanci da baki, kamar yatsu, su ne gaɓɓan hankula
biyar : \VUgras{gani,} \VUgras{ji,} \VUgras{sunsuna,}
\VUgras{ɗanɗano,} \VUgras{taɓawa.}

%%K t02-08-1 p
8. --- Kai fa yana ɗaya daga cikin sassan jikin ɗan adam mafi ban
sha’awa.

%%K t02-tit-2 sub
\VUsaut{4.4mm}
\VUpk{10.2pt}{40}{Gaɓar jiki.}
\VUsaut{3.0mm}

%%K t02-09-1 p
9. --- \VUgras{Wuya} tana haɗa kai da jiki. Tana ɗauke da
\VUgras{maƙogwaro} da \VUgras{ƙeya.}

%%K t02-10-1 p
10. --- Jiki yana kuma ɗauke da gaɓɓai masu muhimmanci da yawa. A
cikin \VUgras{ƙirji} akwai \VUgras{huhu,} waɗanda muke numfashi da su,
da \VUgras{zuciya,} wadda take tsakiyar rai. Idan zuciya ta daina bugu,
mai rai yakan mutu.

%%K t02-11-1 p
11. --- \VUgras{Haƙarƙari} suna riƙe da ƙirji.

%%K t02-12-1 p
12. --- Sauran sassan jiki su ne \VUgras{gefe,} \VUgras{baya,}
\VUgras{kafaɗu} da \VUgras{ciki.} Muna kwanciya a kan gefe ko a kan
baya domin barci, muna kuma ɗaukar kaya a kan kafaɗunmu.

%%K t02-13-1 p
13. --- A cikin ciki akwai \VUgras{tumbi} da \VUgras{hanji.} Suna
taimaka mana mu narkar da abinci.

%%K t02-tit-3 sub
\VUsaut{4.4mm}
\VUpk{10.2pt}{40}{Hannaye da ƙafafu.}
\VUsaut{3.0mm}

%%K t02-14-1 p
14. --- Muna da \VUgras{gaɓɓai} huɗu : \VUgras{hannaye} biyu da
\VUgras{ƙafafu} biyu. Da hannaye muke ɗauka, muke riƙewa, muke ɗagawa
muke kuma tara abubuwa ; da ƙafafu muke tsayawa, muke tafiya muke kuma
gudu.

%%K t02-15-1 p
15. --- \VUgras{Kafaɗa} tana haɗa hannu da jiki. \VUgras{Tsoka} mai
ƙarfi ƙwarai, wato tsokar hannu, ita ce take motsa hannun da
\VUgras{gwiwar hannu} take haɗawa da \VUgras{damtse.}
\VUgras{Wuyan hannu} shi ne gaɓar \VUgras{tafin hannu,} wanda yake
ƙarewa da \VUgras{yatsu} da \VUgras{farata.} A kan hannu muna
rarrabe \VUgras{jijiya} (da jijiyar jini) inda \VUgras{jini} yake
gudana.

%%K t02-16-1 p
16. --- Sassan ƙafa su ne : \VUgras{cinya,} \VUgras{gwiwa} da
\VUgras{bayan gwiwa,} \VUgras{dunƙulen ƙafa,} wanda \VUgras{namansa}
yake rufe da \VUgras{fata,} \VUgras{idon sawu} da \VUgras{tafin
ƙafa.} \VUgras{Ƙasusuwan} ƙafa suna da kauri kuma suna da ƙarfi
ƙwarai. A ƙarshen ƙafa akwai \VUgras{yatsun ƙafa.} Sauran sassa su ne
\VUgras{tafin sawu} da \VUgras{diddige.}

%%K t02-17-1 p
17. --- Haka bayanin jikin ɗan adam yake, kusan gaba ɗaya.

%%K t02-17-2 p
\textit{Bayani.} --- \textit{Da taimakon zancen : « ina jin ciwo a...}
\textit{(kai da sauransu) », malami zai iya sake amfani da wannan}
\textit{dukan kalmomi daga fuskar} \VUgras{lafiyar jiki} \textit{mai
kyau ko mara kyau.}

%%K t02-tit-4 sub
\VUsaut{5.0mm}
\VUpk{10.2pt}{40}{Motsa jiki.}
\VUsaut{2.4mm}
\VUcentre{10.2pt}{0}{\textit{(Labari daga bakin ɗaya daga cikin ɗalibai.)}}
\VUsaut{3.0mm}

%%K t02-18-1 p
18. --- Domin mu zama masu sassauƙar jiki, masu saurin motsi kuma
masu lafiya, dole mu horar da ƙafafunmu da hannayenmu. Shi ya sa muke
wasa muke kuma \VUgras{motsa jiki.}

%%K t02-19-1 p
19. --- A cikin mako, waɗannan horo biyu suna gudana a filin
makarantar sakandare (duba hoto na 1). Amma ranar Alhamis da ranar
Lahadi mukan tafi babban filin ciyawa, inda muke da wuri mu yi gudu mu
kuma nishaɗantar da kanmu.

%%K t02-20-1 p
20. --- Malamanmu da iyayenmu wani lokaci sukan raka mu can ; amma
\VUgras{malamin motsa jiki}\textsuperscript{(12)} shi ne yake jagorantar
horon.

%%K t02-21-1 p
21. --- A kan filin ciyawa, a hannun hagu na ƙofar shiga, akwai
\VUgras{kayan motsa jiki}\textsuperscript{(1)} ; muna hawa
\VUgras{tsani}\textsuperscript{(2)}, muna juyawa a kan
\VUgras{zoban}\textsuperscript{(3)} da a kan \VUgras{sandar
rataya}\textsuperscript{(4)}, muna ɗaga kanmu da hannaye har saman
\VUgras{tsanin igiya}\textsuperscript{(5)} ko na \VUgras{igiya mai
ƙulle-ƙulle}\textsuperscript{(6)}.

%%K t02-22-1 p
22. --- A halin yanzu, abokanmu suna tsallakawa a kan \VUgras{dokin
katako}\textsuperscript{(7)} ko \VUgras{sandunan
kwance}\textsuperscript{(11)}. Waɗansu suna \VUgras{dambe}\textsuperscript{(8)}
ko \VUgras{takobin wasa}\textsuperscript{(9)} da abin rufe fuska da
\VUgras{sandar ƙarfe.} A ƙarshe, waɗansu suna \VUgras{ɗaga
nauyi}\textsuperscript{(10)}.

%%K t02-tit-5 sub
\VUsaut{4.6mm}
\VUpk{10.2pt}{40}{Wasanni.}
\VUsaut{3.0mm}

%%K t02-23-1 p
23. --- Kewaye da masu motsa jiki, yara maza da mata suna wasanni iri
iri. ’Yan mata suna nishaɗantar da kansu da \VUgras{zoben
wasa}\textsuperscript{(14)} ; suna tsalle da \VUgras{igiyar
tsalle}\textsuperscript{(13)} ko su gayyaci ƙannensu maza da ƴan
uwansu zuwa wasan \VUgras{kuroke}\textsuperscript{(15)}. Suna jin daɗin
turar \VUgras{ƙwallon} abokin hamayya da \VUgras{ƙaramin gudumar
katakonsu,} a daidai lokacin da yake tsammani zai ratsa ta ƙarƙashin
baka cikin sauƙi !

%%K t02-24-1 p
24. --- Yara maza sun fi son wasannin da suke buƙatar ƙarfi. Suna son
wasan \VUgras{sandunan wasa}\textsuperscript{(16)} da suke kaɗawa da
\VUgras{ƙwallo}\textsuperscript{(17)} cikin gwaninta. Ba sa kuma rena
wasan \VUgras{abin jujjuyawa}\textsuperscript{(18)} da \VUgras{kama
ƙwallo da sanda}\textsuperscript{(19)} ko ma \VUgras{wasan
kwaɗo}\textsuperscript{(20)}. Amma wasannin da suka fi so su ne
\VUgras{wasan tsakuwa}\textsuperscript{(21)} da \VUgras{ƙwallon
bango}\textsuperscript{(22)}, domin bangon \VUgras{gidan
tsirrai}\textsuperscript{(26)} yana da kyau ƙwarai wajen sake turo
ƙwallon nan maras nauyi.

%%K t02-25-1 p
25. --- Ƙaramin Ioannes, yana tafiya a kan \VUgras{sandunan
ƙafarsa}\textsuperscript{(24)}, gwani ne ƙwarai kuma ya iya kiyaye
daidaito sosai. Kusa da shi, waɗansu abokai suna wasan
\VUgras{ɓoye}\textsuperscript{(25)} a cikin gidan tsirrai da bayan
\VUgras{shinge.} Waɗansu sun fi son \VUgras{zaton mai
bugu}\textsuperscript{(28)} ko \VUgras{ƙwallon
tashi}\textsuperscript{(29)}, wanda yake tashi daga wannan
\VUgras{raket} zuwa wancan.

%%K t02-noto-1 noto
\VUnotes{143.2mm}{%
(*) Wasannin yara sau da yawa suna da sunayen ƙasa-ƙasa. Mun yi ƙoƙari
mu ba su suna a cikin Ido yadda ya fi kyau — wani lokaci muna ɗaukar
tunani daga aikin da yake bayyana su, wani lokaci kuma muna zaɓar
sunan ƙasar nan ko ta can, idan bai yi nisa da gaskiya ba. Misali :
\VUgras{wasan ganga} (na Jamusanci da na Faransanci) shi ne
\VUgras{wasan kwaɗo} \textsuperscript{(20)} (na Turanci), inda masu
wasa suke ƙoƙarin jefa faifansu zagaye ta bakin kwaɗon ƙarfe wanda
yake tsaye a buɗe baki a kan akwatin nan mai ramuka.
\VUgras{Tsallen kafaɗa}\textsuperscript{(30)} ya bambanta da
\VUgras{tsallen baya} a wannan kaɗai : domin ya tsallaka bisa kai,
mai wasa yakan ɗora hannuwansa a kan kafaɗun abokinsa, alhali a
« \VUgras{tsallen baya} » yakan ɗora hannuwansa a bayansa kaɗai. ---
Ko kuma waɗansu daga cikin masu wasa sukan sunkuya alhali waɗansu za
su tsallaka bisa bayansu da hannu a kafaɗa ; na farkon sai su jure
buguwa ba tare da sun tanƙwara ba, na biyun kuwa sai su yi tsalle ba
tare da sun rasa daidaito ba (Duba hoton kansa).%
}

%%K t02-26-1 p
26. --- A tsakiyar filin ciyawa akwai bishiyoyi biyu. A gindin ɗayansu,
yara maza bakwai ko takwas sun shirya wasan \VUgras{tsallen
kafaɗa}\textsuperscript{(30)} (*). Ba a kowace ƙasa ake sanin wannan
wasa ba ; amma kowa ya san abin da \VUgras{tsallen baya} yake, wanda
yaran ba su yi wasa da shi ba a wannan karo.

%%K t02-tit-6 sub
\VUsaut{5.0mm}
\VUpk{10.2pt}{40}{Wasanni a sararin iska.}
\VUsaut{3.4mm}

%%K t02-27-1 p
27. \VUgras{Wasan makaho}\textsuperscript{(31)} ma yana da daɗi ƙwarai ;
ƴan mata da yara maza suna sa \VUgras{ɗaurin ido} a idanunsu bi da
bi suna kama waɗanda ba su iya tsere ba.

%%K t02-28-1 p
28. --- A tsakiyar filin ciyawa, ga wasa mai Faransanci ƙwarai amma
mai ɗan tsanani : shi ne \VUgras{wasan beyar}\textsuperscript{(32)},
mai hayaniya fiye da wasan \VUgras{tsuntsun
iska}\textsuperscript{(33)} mai natsuwar nan, ko ma fiye da wasan
Turancin nan \VUgras{tanis}\textsuperscript{(34)}.

%%K t02-29-1 p
29. --- Wannan wasa na ƙarshe shi ne farin cikin manyan ɗalibai.
Malamanmu da kansu suna yin sa da so. Yana buƙatar gwaninta da saurin
motsi, ni kuwa ina son sa ƙwarai. Na bar wa ƴan mata wasan
\VUgras{lilo}\textsuperscript{(35)}.

%%K t02-30-1 p
30. --- Ba na son wani wasan Turanci wanda wataƙila zai zama mai
haɗari.

%%K t02-30-1b p
Ina nufin \VUgras{ƙwallon ƙafa}\textsuperscript{(36)}, wanda yake faranta
wa babban ɗan uwana rai. Na fi son tsohon \VUgras{wasan kamun
juna}\textsuperscript{(38)} namu, wanda yake da kyau ga lafiya haka nan
kuma ba shi da tsanani sosai.

%%K t02-tit-7 sub
\VUsaut{4.6mm}
\VUpk{10.2pt}{40}{Hawan keke da wasan ƙwallo.}
\VUsaut{3.0mm}

%%K t02-31-1 p
31. --- Ni ma ina son zagaya filin ciyawa a kan
\VUgras{keke}\textsuperscript{(37)}.

%%K t02-32-1 p
32. --- A shekara mai zuwa zan koyi \VUgras{hawan
doki}\textsuperscript{(39)}. Doki yana da kyau ƙwarai idan yana takawa
ko yana sukuwa cikin annashuwa, kuma idan yana tsallake shamaki da
gudu !

%%K t02-33-1 p
33. --- A lokacin rani muna koyon \VUgras{iyo}\textsuperscript{(40)}.
Muna nutsawa muna kuma yin wanka cikin ruwan kogi mai tsafta da sanyi
da daɗi. Wani lokaci a ƙarshe mukan saki \VUgras{balon
takarda}\textsuperscript{(41)} wanda yake hawa sama cikin girma, da
zarar an cika shi da iska mai zafi.

%%K t02-34-1 p
34. --- Akwai kuma wasanni da yawa waɗansu, amma ba zan gama ƙidaya su
ba, kuma daga wannan taƙaitawa mutum yakan gane cewa a ranakun Alhamis
ba a gundura sosai a kan kyakkyawan filin ciyawarmu.

%%K t02-34-2 p
N.-B. --- \textit{Malami zai iya cika wannan babi cikin sauƙi ; ta}
\textit{wurin bayyana kowane ɗaya daga cikin wasannin da suka fi}
\textit{muhimmanci daki-daki.}
